\section*{第 \ref{ch:g11:circuits-and-power} 章——电路与功率}

\begin{solution}{exo:g11:circuits-and-power:1}
$q = It = 2.0 \times 90 \times 60 = \qty{1.1e4}{C}$。元电荷
个数为 $q/e = 1.08\times10^{4}/1.6\times10^{-19} \approx \num{6.8e22}$。
\end{solution}

\begin{solution}{exo:g11:circuits-and-power:2}
$I = U/R = 12/220 = \qty{0.055}{A} = \qty{55}{mA}$;
$R = U/I = 230/0.50 = \qty{460}{\ohm}$.
\end{solution}

\begin{solution}{exo:g11:circuits-and-power:3}
串联：$120 + 60 = \qty{180}{\ohm}$。并联：$\frac{120 \times 60}{180}
= \qty{40}{\ohm}$——比两者都小：电压相同，但电流多了一条
通路，因而总电流更大，电阻更小。
\end{solution}

\begin{solution}{exo:g11:circuits-and-power:4}
$P = UI = 230 \times 8.7 = \qty{2.0}{kW}$。当 $t = \qty{240}{s}$ 时：
$E = Pt = 2000 \times 240 = \qty{4.8e5}{J} = 4.8\times10^{5}/3.6
\times10^{6} \approx \qty{0.13}{kW.h}$.
\end{solution}

\begin{solution}{exo:g11:circuits-and-power:5}
$E = 60 \times 5.0 = \qty{300}{W.h} = \qty{0.30}{kW.h} = 0.30 \times
3.6\times10^{6} = \qty{1.1e6}{J}$。费用：$0.30 \times 0.25 = 0.075$
欧元/天。
\end{solution}

\begin{solution}{exo:g11:circuits-and-power:6}
\emph{1.} $R = U^2/P = 230^2/1200 = \qty{44}{\ohm}$;
$I = P/U = 1200/230 = \qty{5.2}{A}$.

\emph{2.} $P' = U'^2/R = 115^2/44.1 = \qty{300}{W}$：功率变为原来的
四分之一而非一半，因为减半 $U$ 时 $I$ 也随之减半，而 $P = UI \propto U^2$。
\end{solution}

\begin{solution}{exo:g11:circuits-and-power:7}
$R_{\text{串}} = \qty{60}{\ohm}$，故 $I = 12/60 = \qty{0.20}{A}$。
$U_1 = R_1 I = \qty{8.0}{V}$，$U_2 = \qty{4.0}{V}$（总和 \qty{12}{V}）；
$P_1 = U_1 I = \qty{1.6}{W}$，$P_2 = \qty{0.8}{W}$：总功率 \qty{2.4}{W}
$= UI = 12 \times 0.20$，与此相符。
\end{solution}

\begin{solution}{exo:g11:circuits-and-power:8}
$I_1 = 12/60 = \qty{0.20}{A}$，$I_2 = 12/30 = \qty{0.40}{A}$，总电流
$I = \qty{0.60}{A}$。等效电阻：$R = U/I = 12/0.60 = \qty{20}{\ohm}$，
且确实 $\frac{60 \times 30}{60 + 30} = \qty{20}{\ohm}$。
\end{solution}

\begin{solution}{exo:g11:circuits-and-power:9}
$\mathcal{E} - 0.50\,r = 5.5$ 与 $\mathcal{E} - 2.0\,r = 4.0$；
两式相减，$1.5\,r = 1.5$，故 $r = \qty{1.0}{\ohm}$，$\mathcal{E} =
\qty{6.0}{V}$。短路时：$\mathcal{E}/r = \qty{6.0}{A}$。
\end{solution}

\begin{solution}{exo:g11:circuits-and-power:10}
电水壶 $2200/230 = \qty{9.6}{A}$，烤面包机 $1000/230 = \qty{4.3}{A}$：
合计 \qty{13.9}{A} $< \qty{16}{A}$，保险丝不熔断。取暖器 $1500/230 =
\qty{6.5}{A}$：合计 \qty{20.4}{A} $> \qty{16}{A}$，保险丝熔断
（等价地，$\qty{4.7}{kW} > 230 \times 16 \approx \qty{3.7}{kW}$）。
\end{solution}

\begin{solution}{exo:g11:circuits-and-power:11}
在 \qty{200}{V} 下：$I = 10^4/200 = \qty{50}{A}$，损耗 $RI^2 = 2.0 \times
50^2 = \qty{5.0}{kW}$——占输送功率的一半。在 \qty{4.0}{kV} 下：$I =
\qty{2.5}{A}$，损耗 $2.0 \times 2.5^2 = \qty{12.5}{W}$。比值 $400 =
20^2$：电压变为原来的二十倍，意味着电流变为原来的二十分之一，
而损耗与 $I^2$ 成正比。
\end{solution}

\begin{solution}{exo:g11:circuits-and-power:12}
四种连接方式：全部串联，$3 \times 60 = \qty{180}{\ohm}$；全部
并联，$60/3 = \qty{20}{\ohm}$；一个与两个并联的电阻串联，
$60 + 30 = \qty{90}{\ohm}$；一个与两个串联的电阻并联，
$\frac{60 \times 120}{180} = \qty{40}{\ohm}$.
\end{solution}

\begin{solution}{exo:g11:circuits-and-power:13}
\emph{1.} $I = \mathcal{E}/(R + r) = 4.5/7.5 = \qty{0.60}{A}$;
$U = RI = \qty{3.6}{V}$；有用功率 $P = UI = \qty{2.2}{W}$；损失的功率 $rI^2 =
1.5 \times 0.60^2 = \qty{0.54}{W}$; $\eta = U/\mathcal{E} = 3.6/4.5 =
0.80$.

\emph{2.} 由提示，$P = \mathcal{E}^2 \big/ \big[(R - r)^2/R +
4r\big]$；括号项最小（$= 4r$）恰在 $R = r$ 时取到。
$P_{\max} = \mathcal{E}^2/(4r) = 4.5^2/6.0 = \qty{3.4}{W}$.
\end{solution}

\begin{solution}{exo:g11:circuits-and-power:14}
热量：$Q = 1.0 \times 4180 \times (100 - 20) = \qty{3.3e5}{J}$。
电能：$E = Q/0.90 = \qty{3.7e5}{J}$。时间：$t = E/P =
3.7\times10^{5}/2200 \approx \qty{170}{s}$，不到三分钟。费用：
$3.7\times10^{5}/3.6\times10^{6} \approx \qty{0.10}{kW.h}$，约 $2.6$
美分。
\end{solution}

\begin{solution}{exo:g11:circuits-and-power:15}
\emph{1.} $q = 60 \times 3600 = \qty{2.2e5}{C}$; $E \approx \mathcal{E}q
= 12.7 \times 2.16\times10^{5} = \qty{2.7e6}{J} \approx \qty{0.76}{kW.h}$.

\emph{2.} $U = 12.7 - 0.020 \times 150 = \qty{9.7}{V}$；输出功率
$UI = \qty{1.5}{kW}$；损失的功率 $rI^2 = 0.020 \times 150^2 = \qty{450}{W}$。

\emph{3.} 车头灯并联在接线柱两端：启动发动机时，
端电压由 $12.7$ 降至 $\qty{9.7}{V}$，因此车头灯获得的
功率减小，灯光变暗。
\end{solution}

\begin{solution}{pb:g11:circuits-and-power:1}
\textbf{1.} 每秒，$E_p = mgh = 5.0\times10^{4} \times 9.81 \times
200 = \qty{9.8e7}{J}$：可用功率 \qty{98}{MW}。

\textbf{2.} $0.90 \times 98 = \qty{88}{MW}$。

\textbf{3.} $I = P/U = 8.8\times10^{7}/2.0\times10^{4} = \qty{4.4e3}{A}$。

\textbf{4.} $R_\ell I^2 = 32 \times (4.4\times10^{3})^2 \approx
\qty{6.2e8}{W} = \qty{620}{MW}$——是电站输出功率的七倍：
在 \qty{20}{kV} 电压下输电是不可能的。

\textbf{5.} 在 $P = UI$ 一定时，只有提高 $U$ 才能使 $I$ 减小：需要
一台变压器。

\textbf{6.} $I = 8.8\times10^{7}/4.0\times10^{5} = \qty{220}{A}$。

\textbf{7.} $R_\ell I^2 = 32 \times 220^2 \approx \qty{1.6}{MW}$。

\textbf{8.} $1.6/88 \approx 0.018$：约 $1.8\%$。

\textbf{9.} 比值 $= R_\ell I^2/P = R_\ell (P/U)^2/P = R_\ell P/U^2
\propto 1/U^2$：从 \qty{20}{kV} 升至 \qty{400}{kV}，比值减小到原来的
$20^2 = 400$ 分之一；事实上 $\qty{620}{MW}/400 \approx \qty{1.6}{MW}$。

\textbf{10.} $R_\ell I = 32 \times 220 \approx \qty{7.1}{kV}$：占
\qty{400}{kV} 的 $1.8\%$，与第 8 题比例相同，因为
$R_\ell I/U = R_\ell I^2/UI$。

\textbf{11.} $\qty{1}{kW.h} = \qty{3.6e6}{J}$：电表把用电功率
与用电时间相乘，并将结果累加起来。

\textbf{12.} $I = 900/230 = \qty{3.9}{A}$；
$R = U^2/P = 230^2/900 \approx \qty{59}{\ohm}$。

\textbf{13.} $E = 900 \times 180 = \qty{1.6e5}{J} = \qty{0.045}{kW.h}$：
约 $1.1$ 美分。

\textbf{14.} $RI^2$ 的物理原理相同，只是用途
不同——在输电线上热量是不需要的，在烤面包机中热量
\emph{正是}所需的成果。

\textbf{15.} $3.9 + 9.6 = \qty{13.5}{A} < \qty{16}{A}$：成立。加热器
$2000/230 = \qty{8.7}{A}$：总电流 $\qty{22.2}{A} > \qty{16}{A}$，保险丝
熔断。

\textbf{16.} $0.05 \times 13.5^2 = \qty{9.1}{W}$：$9.1/3100 \approx
0.3\%$ 占输送功率——在短导线上可忽略不计。

\textbf{17.} $\eta = 0.90 \times 0.99 \times 0.982 \times 0.99 \times
0.997 \approx 0.86$。

\textbf{18.} 烤面包所耗的一千瓦时电能，上游需要 $3.6\times10^{6}/0.86
\approx \qty{4.2e6}{J}$：$m = E/(gh) = 4.2\times10^{6}/(9.81
\times 200) \approx \qty{2.1e3}{kg}$，约为 \qty{2.1}{m^3} 的水。

\textbf{19.} $R_\ell P/U^2 = 32 \times 8.8\times10^{7}/(2.0\times
10^{4})^2 \approx 7$。“比值”超过 $1$ 意味着导线
消耗的功率超过它所输送的功率：在该电压下，
输电根本无法实现。

\textbf{20.} 烤面包所耗的一千瓦时电能，使山上约两
吨——\qty{2.1}{m^3}——的水从 \qty{200}{m} 高处落下，并有约
$14\%$ 的能量在途中损耗，主要发生在涡轮机和变压器中。
（\qty{1000}{km} 的输电线上损失不到 $2\%$）。使这一旅程成为可能的
唯一关键在于变压器：升压至 \qty{400}{kV}，
可使电流除以 $20$，焦耳热损耗除以 $400$。
\end{solution}
